\subsection{Surface Density Distributions}

One measure of surface activity is the spatial distribution of molecules in the interfacial region. The density distributions of both \wat~and \suldiox~were calculated for the equilibrium MD simulations. The results presented in Figure \ref{fig:density} show both the water (black) and \suldiox~(red) density distributions across the two simulated slabs. As shown, the single \suldiox~in the neat-water system remained at the water slab surface. The \suldiox~in the saturated system accumulated mostly at the surface, but some residual \suldiox~remained well into the bulk water.

%The density distributions of the \suldiox~were fit to gaussian lineshapes as shown in equation \ref{eq:gaussian} using the common definitions of the variables $A$, $\mu$, and $\sigma$, and given as a function of the distance to the top surface, $z$.

%\begin{equation}
%  \rho(z)=\frac{A}{\sqrt{2 \pi \sigma^2}} e^{-\frac{(z-\mu)^2}{2\sigma^2}}
%  \label{eq:gaussian}
%\end{equation}

In the simulated neat-water slab, the result that the \suldiox~molecular density distribution concentrates near the water surface indicates an affinity for the interfacial region. During the course of the simulation, the \suldiox~does not venture into the bulk water nor does it escape the water surface into the gas phase, but remains located within 5-\angs~of the surface region. This is consistent with what is found experimentally in our laboratory using VSFS,\cite{Tarbuck2005,Tarbuck2006} and also supported by the computational simulations and spectral calculations of Baer et al.\cite{Baer2010} The experimental studies indicated that upon exposure of \suldiox~to a \wat~surface, a layer of solvated \suldiox~forms, modifying the structure of water above the Gibbs dividing surface (GDS). In the experimental system, the \suldiox~bonding interaction with the free OH oscillators is manifested in a red-shift of the free OH frequency indicative of a bonding interaction. 

The saturated solution simulation results indicate that under the high concentration of \suldiox, the \suldiox~accumulation at the surface is increased. However, unlike the neat-water slab, the saturated slab has a non-zero bulk concentration of \suldiox. The added concentration of \suldiox~creates a layer of molecules bound to the top of the water surface. The center of the \suldiox~density distribution is further into the gas phase than for the neat-water surface with a single-\suldiox~molecule. Because of the limitations of the classical model in accurately reproducing the first hydration shell around the \suldiox,\cite{Baer2010} we do not draw specific conclusions about the hydrate complex's geometry. However, the surface affinity of the hydrated \suldiox~is well reproduced by both the classical and ab initio methods.\cite{Baer2010}

\begin{figure}[h!]
	\begin{center}
		\includegraphics[scale=1.0]{images/density/density.png}
		\caption{Molecular density distributions of \wat~(black) and \suldiox~(red) calculated along the long-axis of the simulated cells. The position axis shows the distance from the instantaneous surface location of the \wat~slab, with positive values located above the slab towards the gas phase, and negative values located in the water bulk. Distributions of both the neat-water simulation with a single \suldiox~(left, with the \suldiox~density scaled 10x for clarity), and the saturated system (right) are shown.}
		\label{fig:density}
	\end{center}
\end{figure}
